\section{Verification and provenance}\label{sec:verification}

\paragraph{Exactness.}
A record is a list of triples of bit vectors whose sum is the matrix-multiplication tensor (\cref{sec:setting}). Each of the 795 records was checked by two separately written implementations: a sparse cancellation count in Python and a dense comparison with the closed form in Ruby; the native walker checks a third time the input and output of every walk. Both checkers reject a single flipped bit, a dropped term and an output factor in the transposed layout of the public files (the negative controls are in the companion paper).
The Ruby check ran on the 691 records held in the release bundles and on the 661 chain heads of the bundles (logs in the release), and later on all 1{,}089 tensors behind the records: the 985 distinct tensors of the bundles, the 102 committed certificates that stay in the repository and the two fleet checkpoints (log in the repository); every one is exact at the rank of its manifest or index row, and every bundle file has the hash of its manifest row. The verification runs of the release itself (A--D) used the Python gate only.
Both checkers were written by the same author and share the bit layout, which is the weakest point of this claim.

\paragraph{Identity and lineage.}
A tensor is identified by the SHA-256 of its canonical text (sorted, duplicate-free triples). The index gives, for every record, the format, the rank, the best listed rank in any ring and among GF(2)-valid schemes, the two margins, the class, the lineage family, the SHA-256 and the holding asset. Each record ends a chain of rows, from a seed (a public scheme reduced modulo two, a replayed catalogue recipe, a closure recipe or a committed certificate) through projections, basis rewrites and walks, in one of four bundles.
The verifier re-derived from pinned upstream files or from git history all 11 distinct public seeds, all 186 archived July certificates, matched the 25 retained committed certificates against the repository files, recomputed all 225 projection rows and reproduced 71 of the 82 catalogue replays from an empty cache; and it replayed 42 of the 244 distinct walks byte for byte with the native walker. The other walks were checked for exactness and for the invariant that a walk keeps the format and never raises the rank.
Running the shipped verifier on the extracted release reproduces its five logs exactly.

\paragraph{What is not established.}
Five replayed seeds, and the eight records that depend on them, were built from two local leaves (schemes for $\MM448$ and $\MM{12}{15}{16}$) that their manifest does not list: the tensors are exact, but their leaf provenance is incomplete, and the archive's description of them as ``equal-rank variants'' of an empty-cache replay is wrong (they are lower-rank first-pass constructions).
A separate script rebuilt 259 of the 267 closure rows byte for byte; the other eight rows, six distinct records, depend on three intermediate witnesses whose files no longer exist, and 43 of 1{,}098 replays give a tensor of the same rank that differs from what an empty cache gives, through a cache-order defect that is fixed and described in the companion paper. None of this affects exactness or rank.
The record index also carries two wrong lineage labels, corrected in this paper: it lists $\MM999$ and $\MM{13}{13}{13}$ under ``live fleet checkpoint'', whereas the first is a walk from a public scheme and the second the July certificate itself.
The steps by which the 126 committed certificates were produced before the campaign are documented in the repository's certificate directories, and were not replayed for this release.
